\chapter*{Introduction générale}
\addcontentsline{toc}{chapter}{Introduction générale}
\markboth{Introduction générale}{}

Les véhicules connectés et autonomes reposent sur une idée simple : un véhicule qui perçoit son environnement,
qui communique ce qu'il perçoit et qui en tire des décisions est plus sûr et plus efficace. Cette idée se trouve
au croisement de deux domaines qui progressent rapidement : l'Internet des objets (IoT), qui relie des capteurs et
des actionneurs à des services distants, et l'intelligence artificielle, et plus particulièrement la vision par
ordinateur, qui permet d'interpréter automatiquement des images.

Sur la route, les obstacles ne sont pas seulement les autres usagers. Les nids-de-poule, les cônes, les barrières
et les zones de travaux sont à l'origine de dégâts matériels et d'accidents, et leur signalement reste souvent
manuel et tardif. Un système capable de les détecter automatiquement à partir d'une caméra embarquée, de les
localiser, d'en évaluer la gravité et de les partager avec les autres véhicules apporterait une aide directe à la
conduite et à la maintenance du réseau routier.

C'est dans ce contexte que s'inscrit notre projet, réalisé dans le cadre du module \emph{Internet des objets} à la
Faculté des Sciences et Techniques de Béni Mellal. Nous avons conçu une chaîne complète, du capteur jusqu'à
l'interface de supervision :
\begin{itemize}
  \item un \textbf{prototype} de voiture miniature qui suit une piste et réagit aux obstacles, équipé d'une
        caméra Wi-Fi ESP32-CAM ;
  \item des \textbf{modèles de détection d'objets} YOLO26, entraînés d'une part sur des données routières réelles
        et d'autre part sur des images de notre propre piste ;
  \item une \textbf{application web}, \app, qui exécute ces modèles en temps réel, estime la gravité des
        obstacles, les priorise, les place sur une carte et propose des itinéraires qui les évitent.
\end{itemize}

L'application gère deux environnements. L'environnement \emph{monde réel} simule une petite flotte de véhicules
sur le réseau routier réel d'un quartier de Paris ; la caméra du véhicule CAR-01 y est remplacée par de vraies
vidéos de route analysées par nos modèles. L'environnement \emph{prototype} correspond à la voiture miniature et
à sa piste. Ce double environnement permet de valider la même chaîne logicielle à deux échelles.

Ce rapport est organisé en quatre chapitres. Le \textbf{premier chapitre} présente le contexte général, l'IoT et
son association avec l'intelligence artificielle, la problématique, les objectifs, le cahier des charges et
l'architecture générale du système. Le \textbf{deuxième chapitre} décrit la conception et la réalisation
matérielle du prototype. Le \textbf{troisième chapitre} est consacré aux modèles d'intelligence artificielle :
données, entraînement, évaluation et intégration. Le \textbf{quatrième chapitre} présente l'application
\app{} page par page. Nous terminons par une conclusion générale et des perspectives.
